% Part: history 
% Chapter: biographies 
% Section: julia-robinson
\documentclass[../../../include/open-logic-section]{subfiles}

\begin{document}

\olfileid{his}{bio}{rob}

\olsection{ज्युलिया रॉबिन्सन}

\olphoto{robinson-julia}{ज्युलिया रॉबिन्सन}

ज्युलिया बोमन रॉबिन्सन या अमेरिकन गणितज्ञ होत्या. त्या मुख्यतः निर्णय समस्यांवरील कार्यासाठी, आणि विशेषतः हिल्बर्ट यांच्या दहाव्या समस्येच्या समाधानातील योगदानासाठी ओळखल्या जातात. रॉबिन्सन यांचा जन्म ८ डिसेंबर १९१९ रोजी मिसुरी राज्यातील सेंट लुईस येथे झाला. लहानपणापासूनच संख्यांविषयी कुतूहल होते, अशी त्यांची आठवण आहे \citep[4]{Reid1986}. वयाच्या नवव्या वर्षी त्यांना स्कार्लेट फीव्हर झाला आणि नंतर संधिवातजन्य तापाचे अनेक झटके आले. त्यामुळे त्यांना बराच वेळ अंथरुणात घालवावा लागला आणि शिक्षणात त्या मागे पडल्या. खाजगी शिक्षकांच्या मदतीने त्यांनी ही तूट भरून काढली; मात्र आजाराचे शारीरिक परिणाम त्यांच्या पुढील आयुष्यावरही राहिले.

बालपणीच्या अडचणी असूनही रॉबिन्सन यांनी गणित आणि विज्ञानातील अनेक पुरस्कारांसह माध्यमिक शिक्षण पूर्ण केले. विद्यापीठीय शिक्षणाची सुरुवात त्यांनी सॅन डिएगो स्टेट कॉलेजमध्ये केली आणि शेवटच्या वर्षी कॅलिफोर्निया विद्यापीठ, बर्कली येथे प्रवेश घेतला. तेथे गणितज्ञ राफाएल रॉबिन्सन यांचा त्यांच्यावर प्रभाव पडला. दोघांची मैत्री झाली आणि १९४१ मध्ये त्यांनी विवाह केला. विद्याशाखेतील सदस्याची पत्नी असल्याने ज्युलिया यांना बर्कलीच्या गणित विभागात अध्यापनाची परवानगी नव्हती. त्या गणिताचे वर्ग श्रोती म्हणून घेत राहिल्या, पण विद्यापीठ सोडून कुटुंब सुरू करण्याची त्यांची इच्छा होती. विवाहानंतर थोड्याच काळात त्यांना न्यूमोनिया झाला. बालपणीच्या संधिवातजन्य तापामुळे त्यांच्या हृदयावर बरीच व्रणयुक्त ऊतकांची वाढ झाली आहे, असे त्यांना सांगण्यात आले. या स्थितीची तीव्रता पाहून डॉक्टरांनी त्या चाळीस वर्षांपलीकडे जगणार नाहीत असा अंदाज व्यक्त केला आणि त्यांना मुले न होऊ देण्याचा सल्ला दिला \citep[13]{Reid1986}.

रॉबिन्सन दीर्घकाळ नैराश्यात होत्या; अखेरीस त्यांनी गणिताचे शिक्षण सुरू ठेवायचे ठरवले. त्या बर्कलीला परतल्या आणि ॲल्फ्रेड टार्स्की यांच्या मार्गदर्शनाखाली १९४८ मध्ये डॉक्टरेट पूर्ण केली. वास्तविक संख्यांची प्रथम-क्रम उपपत्ती निर्णेय आहे, हे टार्स्की यांनी दाखवले होते; गोडेल यांच्या कार्यावरून नैसर्गिक संख्यांची प्रथम-क्रम उपपत्ती अनिर्णेय आहे, हे निष्पन्न होत होते. परिमेय संख्यांची प्रथम-क्रम उपपत्ती निर्णेय आहे की नाही, हा मोठा अनिर्णित प्रश्न होता. आपल्या प्रबंधात \citeyearpar{Robinson1949} रॉबिन्सन यांनी ती अनिर्णेय आहे, हे सिद्ध केले.

निर्णय समस्यांत रस असल्याने रॉबिन्सन यांनी पुढे हिल्बर्ट यांच्या दहाव्या समस्येचे समाधान शोधण्याचा प्रयत्न केला. १९०० मध्ये डेव्हिड हिल्बर्ट यांनी मांडलेल्या २३ प्रसिद्ध गणिती समस्यांपैकी ही एक होती. पूर्णांक गुणांक असलेल्या $3x^2 - 2y +3 = 0$ यांसारख्या बहुपदी समीकरणाला पूर्णांकांमध्ये उत्तर आहे की नाही, हे मर्यादित वेळेत सांगणारा कलनविधी आहे का, असा दहाव्या समस्येचा प्रश्न आहे. अशा प्रश्नांना \emph{डायोफँटाइन समस्या} म्हणतात. काही सुरुवातीच्या यशांनंतर रॉबिन्सन यांनी याच समस्येवर काम करणारे मार्टिन डेव्हिस आणि हिलरी पुटनम यांच्यासोबत काम सुरू केले. अज्ञात राशी घातांकांतही येऊ शकतात अशा घातांकीय डायोफँटाइन समस्या अनिर्णेय आहेत, हे त्यांनी दाखवले. तसेच पुढे “J.R.” म्हणून ओळखल्या गेलेल्या एका विशिष्ट गृहीतकावरून हिल्बर्ट यांची दहावी समस्या अनिर्णेय आहे हे निष्पन्न होते, असेही त्यांनी सिद्ध केले \citep{DavisPutnamRobinson1961}. १९६०च्या दशकातही रॉबिन्सन या समस्येवर काम करत राहिल्या. अखेरीस १९७० मध्ये तरुण रशियन गणितज्ञ युरी मातियासेविच यांनी J.R. गृहीतक सिद्ध केले. एकत्रित निष्कर्षाला आता मातियासेविच–रॉबिन्सन–डेव्हिस–पुटनम प्रमेय, संक्षेपात MRDP प्रमेय, म्हणतात. मातियासेविच आणि रॉबिन्सन यांची मैत्री झाली आणि त्यांनी अनेक लेखांवर एकत्र काम केले. मातियासेविच यांना लिहिलेल्या एका पत्रात रॉबिन्सन यांनी असे म्हटले होते: “प्रत्यक्षात, हजारो मैल दूर राहूनही एकत्र काम केल्यामुळे आपल्यापैकी कोणीही एकट्याने जितकी प्रगती केली असती त्याहून अधिक प्रगती आपण करत आहोत, याचा मला खूप आनंद आहे” \citep[45]{Matijasevich1992}.

रॉबिन्सन या अमेरिकन मॅथेमॅटिकल सोसायटीच्या पहिल्या महिला अध्यक्ष होत्या. राष्ट्रीय विज्ञान अकादमीच्या गणित विभागात निवडून जाणाऱ्या त्या पहिल्या महिला होत्या. रक्ताच्या कर्करोगाचे निदान झाल्यानंतर ३० जुलै १९८५ रोजी वयाच्या ६५व्या वर्षी त्यांचे निधन झाले.

\begin{reading}
रॉबिन्सन यांचे गणितीय लेख त्यांच्या \textit{Collected Works} \citep{Robinson1996} मध्ये उपलब्ध आहेत. त्यात राष्ट्रीय विज्ञान अकादमीसाठी लिहिलेल्या त्यांच्या चरित्रात्मक स्मृतिलेखाचे पुनर्मुद्रणही आहे \citep{Feferman1994}. रॉबिन्सन यांची मोठी बहीण कॉन्स्टन्स रीड यांनी मुलाखतींवर आधारित “Autobiography of Julia” \citep{Reid1986} आणि त्यांचे सविस्तर चरित्र \citep{Reid1996} प्रकाशित केले. रॉबिन्सन आणि हिल्बर्ट यांच्या दहाव्या समस्येवरील लघुपटाचे दिग्दर्शन जॉर्ज चिचेरी यांनी केले \citep{Csicsery2016}. मातियासेविच यांचे रॉबिन्सन यांच्यासोबतचे सहकार्य आणि त्यांच्या कार्यावरील रॉबिन्सन यांचा प्रभाव यांविषयीच्या संक्षिप्त स्मृतिलेखासाठी \citep{Matijasevich1992} पाहा.
\end{reading}

\end{document}
